% !TEX root = ../nc-quadric-classification.tex
% 2026-11-21 / v1.34
\section{可換 quadric の全 Artin 変形と相対 AS4 被覆}
\label{sec:quadric-artin-hull}

以下では $(A,\mathfrak m,k)$ を標数零の体上の有限次元可換局所 Artin 環とする。
本節は一般の相対 canonical hull 定理~\ref{thm:universal-cubic-hull} を
可換 quadric の Artin 変形へ適用する。
主眼は、任意の pointed Abel 変形から、同じ構造対象を持つ exact ample な
Serre 平方根と相対 AS4 被覆を回復することである。
二周期の $\mathbb Z$-代数を用い、cubic 代数の一周期性は仮定しない。

\paragraph{入力を有限自由な tensor データとして固定する。}
$D$ を可換 quadric の cubic $\mathbb Z$-代数の $A$-平坦変形とし、
$V_i=D_{i,i+1}$、$R_i\subset V_i\otimes V_{i+1}\otimes V_{i+2}$ を
それぞれ階数二の一次生成元・三次関係の自由加群とする。
三次の標準解消は相対完全であり、
\[
 W_i=(R_i\otimes V_{i+3})\cap(V_i\otimes R_{i+1})
\]
は階数一で、任意の生成元 $w_i$ は両方の表示で非退化とする。
この非退化性は閉ファイバーでの非退化性と Nakayama から従う。
次数二の周期 $\beta:D\to D(2)$ が与えられたとし、
その偶数部分の graded 化を $B$ とする。
特に $B_1$ は階数四、$B_n$ は階数 $(n+1)^2$ である。

$B$ は二次環で相対 Koszul、二次双対 $E=B^!$ の各次数は有限自由で、
その階数は $1,4,7,8,8,\ldots$ と仮定する。
これらは後述の helix からの入力では既存の相対 Koszul 補題により成立する。
双対の積は tensor 部分加群の包含を双対化する規約で書く。
全ての $i$ に対して
\[
 S_i=R_i\otimes V_{i+3}+V_i\otimes R_{i+1}
 \subset V_i\otimes V_{i+1}\otimes V_{i+2}\otimes V_{i+3}
\]
と置く。これは階数七であり、偶数 $i$ では $\check B$ の二次関係加群である。

\begin{lemma}[平坦な cubic 変形の syzygy と四周期]
\label{lem:relative-cubic-tensors-artin}
$D$ の各成分が有限自由で、閉ファイバーが可換 quadric の cubic
$\mathbb Z$-代数であるとする。
このとき上記の三次提示、$R_i,W_i$ の自由性と底変更、非退化性、
三次の標準解消の相対完全性が成立する。
また生成元 $w_i\in W_i$ を選ぶと、$w_i$ の巡回から
$\theta_i:V_i\simeq V_{i+4}$ が定まり、$D\simeq D(4)$ を誘導する。
\end{lemma}
\begin{proof}
証明は相対モジュライ別稿 \cite[\texttt{quadric_artin_hull} 節]{QuadricModuli2026} に置く。
ここでは後続の分類側から参照するため、仮定と主張を保持する。
\end{proof}

\begin{lemma}[canonical functional の相対的な正規性]
\label{lem:relative-canonical-hull-functional}
$w_i$ が定める二つの coevaluation に対する評価写像を
$R_i\otimes V_{i+3}$ と $V_i\otimes R_{i+1}$ 上に取る。
両者は $W_i$ 上で一致し、全射
\[
 \phi_i:S_i\longrightarrow A,\qquad \phi_i(w_i)=2
\]
を定める。その核 $H_i$ は階数六の直和因子である。
$\phi_i$ の張る直線は $w_i$ の単元倍によらず、$\beta$ と整合する。
さらに $E$ の二次元に対応するこれらの元は、正規化列をなす。
従って $\beta$ に沿って graded 化すれば、$E_2$ に正規元
$\phi$ が得られる。
\end{lemma}
\begin{proof}
証明は相対モジュライ別稿 \cite[\texttt{quadric_artin_hull} 節]{QuadricModuli2026} に置く。
ここでは後続の分類側から参照するため、仮定と主張を保持する。
\end{proof}

\begin{proposition}[正規 functional からの相対 AS4 被覆]
\label{prop:relative-artin-hull-from-functional}
前補題の入力を満たすとき、$B$ は全次数平坦な二次 AS4 被覆
$\widetilde B\twoheadrightarrow B$ を持ち、その核は二次の左右正則正規元で生成される。
\end{proposition}
\begin{proof}
証明は相対モジュライ別稿 \cite[\texttt{quadric_artin_hull} 節]{QuadricModuli2026} に置く。
ここでは後続の分類側から参照するため、仮定と主張を保持する。
\end{proof}

\begin{remark}[可換ファイバーでの正規化]
標準二周期を持つ Segre 環では、$\phi_0$ は六つの交換関係を零に送り、
$q=x_{00}x_{11}-x_{01}x_{10}$ を $1$ に送る。
これは \cite[(5.3)]{VdB2011} の
\[
 w=u_0v_0u_1v_1-u_1v_0u_0v_1-u_0v_1u_1v_0+u_1v_1u_0v_0
\]
を逆 coevaluation の定義へ代入して確かめられる。
従って $E_k/(\phi_0)$ は四変数の外積代数であり、
その双対 hull は四変数多項式環となる。
\end{remark}

\begin{theorem}[可換 quadric の全 Artin 変形の相対実現]
\label{thm:quadric-all-artin-h2}
$Q=\mathbb P^1_k\times\mathbb P^1_k$ とし、
$(\mathcal C_A,\mathcal O_A)$ を $(\operatorname{coh}Q,\mathcal O_Q)$ の
任意の flat pointed Abel 変形とする。
この変形には $s_0=(-)\otimes\mathcal O_Q(1,1)$ を還元とする
exact ample な相対 Serre 平方根 $s_A$ が存在し、その全軌道環 $B_A$ は
\[
 B_A\simeq\widetilde B_A/(g),\qquad
 \operatorname{rank}_A(B_A)_n=(n+1)^2
\]
と表示できる。ここで $\widetilde B_A$ は型 $(1,4,6,4,1)$ の相対 quadratic
AS 正則環であり、$g$ は二次の左右正則正規元である。
従って $(\mathcal C_A,\mathcal O_A)$ は、任意の有限型 $A$-代数への基底変更を
備えた定義~\ref{def:relative-h2} の相対 $\mathfrak H_2$ 族へ昇格する。
根と被覆は存在条件であり、新しい標識として加えない。
\end{theorem}
\begin{proof}
証明は相対モジュライ別稿 \cite[\texttt{quadric_artin_hull} 節]{QuadricModuli2026} に置く。
ここでは後続の分類側から参照するため、仮定と主張を保持する。
\end{proof}

\begin{corollary}[可換点での形式変形問題]
\label{cor:quadric-formal-deformation-equivalence}
$Q$ の相対 $\mathfrak H_2$ 変形から flat pointed Abel 変形への忘却は、
局所 Artin $k$-代数上の変形 groupoid の同値を与える。
この変形問題は障害を持たず、その接空間は
\[
 \operatorname{HH}^2(Q)=H^0(Q,\mathcal O_Q(2,2)),\qquad \dim=9
\]
である。無限小自己同型の空間は $\operatorname{HH}^1(Q)=H^0(Q,T_Q)$、次元六である。
\end{corollary}
\begin{proof}
証明は相対モジュライ別稿 \cite[\texttt{quadric_artin_hull} 節]{QuadricModuli2026} に置く。
ここでは後続の分類側から参照するため、仮定と主張を保持する。
\end{proof}

この定理は可換点における全 Artin 基底の認識を与える。
一般有限型基底への代数化、一般の非可換点での同様の認識、
内在的 $\mathfrak H_2$ 全体の代数性は別に証明する必要がある。

\subsection{cubic 点での条件付き Artin 認識}

\begin{proposition}[適合 cubic helix と Hochschild 自明作用からの全 Artin 認識]
\label{prop:cubic-artin-hhtrivial-conditional}
標数零の代数閉体 $k$ 上の $(\mathcal C,\mathcal O)\in\mathcal H_2$ について、
次のデータと条件を仮定する。
\begin{enumerate}
\item $E_0=\mathcal O$ を含む ample geometric four-helix $(E_i)_{i\in\mathbb Z}$ があり、
  $D_{ij}=\operatorname{Hom}(E_{-j},E_{-i})$ は標準的な三次三次元 AS 正則
  $\mathbb Z$-代数である。
\item exact な Serre 平方根 $s_0^2\simeq S^{-1}[2]$ があり、
  $(\mathcal O,s_0)$ は Artin--Zhang ample で、
  全ての $i$ について $s_0E_i\simeq E_{i+2}$ である。
  ここでは対象ごとの同型で十分であり、別途 coherent な周期の標識は要求しない。
\item $(s_0)_*$ は $\operatorname{HH}^2(\mathcal C)$ 上恒等に作用する。
\end{enumerate}
このとき、任意の局所 Artin $k$-代数 $A$ 上の flat pointed Abel 変形
$(\mathcal C_A,\mathcal O_A)$ は、$s_0$ の持上げを用いる全軌道環の
相対 quadratic AS4 被覆を持つ。
従って全有限型基底変更を含む相対 $\mathfrak H_2$ 族へ昇格する。
この点での相対 $\mathfrak H_2$ 変形と flat pointed Abel 変形の groupoid は、
局所 Artin 基底上で同値となる。

さらに $(s_0)_*$ が $\operatorname{HH}^1(\mathcal C)$ 上も恒等なら、
この還元を持つ Serre 平方根の持上げは Morita 自己同値の同型類として一意である。
\end{proposition}
\begin{proof}
証明は相対モジュライ別稿 \cite[\texttt{quadric_artin_hull} 節]{QuadricModuli2026} に置く。
ここでは後続の分類側から参照するため、仮定と主張を保持する。
\end{proof}

\begin{remark}[Hochschild 作用条件は対象ごとの周期性から従わない]
\label{rem:cubic-helix-root-action}
この命題は、全ての cubic 点または全ての平方根について
$(s_0)_*=1$ on $\operatorname{HH}^2$ と主張していない。
$s_0^2=S^{-1}[2]$ から得られるのは作用が対合であることまでである。
実際、可換 $Q$ 上で $\varphi(z,w)=(-z,w)$ とすると
$s_0=(-\otimes\mathcal O_Q(1,1))\circ\varphi^*$ は Serre 平方根であり、
標準 helix の各線束を対象ごとに二項先へ送る。
しかし $\partial_z\wedge\partial_w$ を負に送るので、
$\operatorname{HH}^2$ 上恒等ではない。
従って cubic helix との対象ごとの適合性から作用条件を除くことはできない。
一般 cubic 点に前命題を適用するには、適合する根で Hochschild 二次に
恒等に作用するものが存在する範囲を別に確定する必要がある。
\end{remark}
